% Options for packages loaded elsewhere
\PassOptionsToPackage{unicode}{hyperref}
\PassOptionsToPackage{hyphens}{url}
\PassOptionsToPackage{dvipsnames,svgnames,x11names}{xcolor}
%
\documentclass[
  ignorenonframetext,
  aspectratio=169,
]{beamer}
\usepackage{pgfpages}
\setbeamertemplate{caption}[numbered]
\setbeamertemplate{caption label separator}{: }
\setbeamercolor{caption name}{fg=normal text.fg}
\beamertemplatenavigationsymbolsempty
% Prevent slide breaks in the middle of a paragraph
\widowpenalties 1 10000
\raggedbottom
\setbeamertemplate{part page}{
  \centering
  \begin{beamercolorbox}[sep=16pt,center]{part title}
    \usebeamerfont{part title}\insertpart\par
  \end{beamercolorbox}
}
\setbeamertemplate{section page}{
  \centering
  \begin{beamercolorbox}[sep=12pt,center]{part title}
    \usebeamerfont{section title}\insertsection\par
  \end{beamercolorbox}
}
\setbeamertemplate{subsection page}{
  \centering
  \begin{beamercolorbox}[sep=8pt,center]{part title}
    \usebeamerfont{subsection title}\insertsubsection\par
  \end{beamercolorbox}
}
\AtBeginPart{
  \frame{\partpage}
}
\AtBeginSection{
  \ifbibliography
  \else
    \frame{\sectionpage}
  \fi
}
\AtBeginSubsection{
  \frame{\subsectionpage}
}
\usepackage{amsmath,amssymb}
\usepackage{iftex}
\ifPDFTeX
  \usepackage[T1]{fontenc}
  \usepackage[utf8]{inputenc}
  \usepackage{textcomp} % provide euro and other symbols
\else % if luatex or xetex
  \usepackage{unicode-math} % this also loads fontspec
  \defaultfontfeatures{Scale=MatchLowercase}
  \defaultfontfeatures[\rmfamily]{Ligatures=TeX,Scale=1}
\fi
\usepackage{lmodern}
\usetheme[]{metropolis}
\usefonttheme{professionalfonts}
\ifPDFTeX\else
  % xetex/luatex font selection
  \setmonofont[Scale=0.78]{DejaVu Sans Mono}
\fi
% Use upquote if available, for straight quotes in verbatim environments
\IfFileExists{upquote.sty}{\usepackage{upquote}}{}
\IfFileExists{microtype.sty}{% use microtype if available
  \usepackage[]{microtype}
  \UseMicrotypeSet[protrusion]{basicmath} % disable protrusion for tt fonts
}{}
\makeatletter
\@ifundefined{KOMAClassName}{% if non-KOMA class
  \IfFileExists{parskip.sty}{%
    \usepackage{parskip}
  }{% else
    \setlength{\parindent}{0pt}
    \setlength{\parskip}{6pt plus 2pt minus 1pt}}
}{% if KOMA class
  \KOMAoptions{parskip=half}}
\makeatother
\usepackage{xcolor}
\newif\ifbibliography
\usepackage{color}
\usepackage{fancyvrb}
\newcommand{\VerbBar}{|}
\newcommand{\VERB}{\Verb[commandchars=\\\{\}]}
\DefineVerbatimEnvironment{Highlighting}{Verbatim}{commandchars=\\\{\}}
% Add ',fontsize=\small' for more characters per line
\usepackage{framed}
\definecolor{shadecolor}{RGB}{35,38,41}
\newenvironment{Shaded}{\begin{snugshade}}{\end{snugshade}}
\newcommand{\AlertTok}[1]{\textcolor[rgb]{0.58,0.85,0.30}{\textbf{\colorbox[rgb]{0.30,0.12,0.14}{#1}}}}
\newcommand{\AnnotationTok}[1]{\textcolor[rgb]{0.25,0.50,0.35}{#1}}
\newcommand{\AttributeTok}[1]{\textcolor[rgb]{0.16,0.50,0.73}{#1}}
\newcommand{\BaseNTok}[1]{\textcolor[rgb]{0.96,0.45,0.00}{#1}}
\newcommand{\BuiltInTok}[1]{\textcolor[rgb]{0.50,0.55,0.55}{#1}}
\newcommand{\CharTok}[1]{\textcolor[rgb]{0.24,0.68,0.91}{#1}}
\newcommand{\CommentTok}[1]{\textcolor[rgb]{0.48,0.49,0.49}{#1}}
\newcommand{\CommentVarTok}[1]{\textcolor[rgb]{0.50,0.55,0.55}{#1}}
\newcommand{\ConstantTok}[1]{\textcolor[rgb]{0.15,0.68,0.68}{\textbf{#1}}}
\newcommand{\ControlFlowTok}[1]{\textcolor[rgb]{0.99,0.74,0.29}{\textbf{#1}}}
\newcommand{\DataTypeTok}[1]{\textcolor[rgb]{0.16,0.50,0.73}{#1}}
\newcommand{\DecValTok}[1]{\textcolor[rgb]{0.96,0.45,0.00}{#1}}
\newcommand{\DocumentationTok}[1]{\textcolor[rgb]{0.64,0.20,0.25}{#1}}
\newcommand{\ErrorTok}[1]{\textcolor[rgb]{0.85,0.27,0.33}{\underline{#1}}}
\newcommand{\ExtensionTok}[1]{\textcolor[rgb]{0.00,0.60,1.00}{\textbf{#1}}}
\newcommand{\FloatTok}[1]{\textcolor[rgb]{0.96,0.45,0.00}{#1}}
\newcommand{\FunctionTok}[1]{\textcolor[rgb]{0.56,0.27,0.68}{#1}}
\newcommand{\ImportTok}[1]{\textcolor[rgb]{0.15,0.68,0.38}{#1}}
\newcommand{\InformationTok}[1]{\textcolor[rgb]{0.77,0.36,0.00}{#1}}
\newcommand{\KeywordTok}[1]{\textcolor[rgb]{0.81,0.81,0.76}{\textbf{#1}}}
\newcommand{\NormalTok}[1]{\textcolor[rgb]{0.81,0.81,0.76}{#1}}
\newcommand{\OperatorTok}[1]{\textcolor[rgb]{0.81,0.81,0.76}{#1}}
\newcommand{\OtherTok}[1]{\textcolor[rgb]{0.15,0.68,0.38}{#1}}
\newcommand{\PreprocessorTok}[1]{\textcolor[rgb]{0.15,0.68,0.38}{#1}}
\newcommand{\RegionMarkerTok}[1]{\textcolor[rgb]{0.16,0.50,0.73}{\colorbox[rgb]{0.08,0.19,0.26}{#1}}}
\newcommand{\SpecialCharTok}[1]{\textcolor[rgb]{0.24,0.68,0.91}{#1}}
\newcommand{\SpecialStringTok}[1]{\textcolor[rgb]{0.85,0.27,0.33}{#1}}
\newcommand{\StringTok}[1]{\textcolor[rgb]{0.96,0.31,0.31}{#1}}
\newcommand{\VariableTok}[1]{\textcolor[rgb]{0.15,0.68,0.68}{#1}}
\newcommand{\VerbatimStringTok}[1]{\textcolor[rgb]{0.85,0.27,0.33}{#1}}
\newcommand{\WarningTok}[1]{\textcolor[rgb]{0.85,0.27,0.33}{#1}}
\usepackage{longtable,booktabs,array}
\usepackage{calc} % for calculating minipage widths
\usepackage{caption}
% Make caption package work with longtable
\makeatletter
\def\fnum@table{\tablename~\thetable}
\makeatother
\setlength{\emergencystretch}{3em} % prevent overfull lines
\providecommand{\tightlist}{%
  \setlength{\itemsep}{0pt}\setlength{\parskip}{0pt}}
\setcounter{secnumdepth}{-\maxdimen} % remove section numbering
\metroset{block=fill}
\setbeamertemplate{navigation symbols}{}
\usepackage{booktabs}
\usepackage{amsmath,amssymb}
\ifLuaTeX
  \usepackage{selnolig}  % disable illegal ligatures
\fi
\usepackage{bookmark}
\IfFileExists{xurl.sty}{\usepackage{xurl}}{} % add URL line breaks if available
\urlstyle{same}
\hypersetup{
  pdftitle={Lean Agent Code \& Multi-Agent Systems},
  pdfauthor={MQS · mark@mqs.dk},
  colorlinks=true,
  linkcolor={Maroon},
  filecolor={Maroon},
  citecolor={Blue},
  urlcolor={Blue},
  pdfcreator={LaTeX via pandoc}}

\title{Lean Agent Code \& Multi-Agent Systems}
\subtitle{Ponytail · gitea-mcp · Claude Code · Self-Hosted LLMs}
\author{MQS · mark@mqs.dk}
\date{2026}
\institute{Modular Quantum Solutions}

\begin{document}
\frame{\titlepage}

\section{Lean Agent Code}\label{lean-agent-code}

\begin{frame}{Why Code Quality Matters for AI Agents}
\phantomsection\label{why-code-quality-matters-for-ai-agents}
\textbf{The problem}: AI agents tend to over-engineer --- adding layers,
abstractions, and boilerplate that create maintenance debt without
adding value.

\textbf{Research findings on AGENTS.md / code-guidance files}:

\begin{longtable}[]{@{}ll@{}}
\toprule\noalign{}
Intervention & Compliance \\
\midrule\noalign{}
\endhead
Documentation alone & \textasciitilde25--40\% \\
Runtime enforcement hooks & \textasciitilde95\% \\
Auto-generated AGENTS.md & \emph{reduces} success rate by
\textasciitilde3\% \\
\bottomrule\noalign{}
\end{longtable}

\vspace{0.3em}
\begin{block}{Conclusion}
Rules without enforcement drift. The solution is a decision ladder baked into every coding session.
\end{block}
\end{frame}

\begin{frame}{Ponytail Decision Ladder}
\phantomsection\label{ponytail-decision-ladder}
Before writing a single line of code, an AI agent asks:

\begin{enumerate}
\tightlist
\item
  \textbf{Does it need to exist?} --- skip if not required
\item
  \textbf{Already in codebase?} --- reuse existing code
\item
  \textbf{Stdlib handles it?} --- prefer standard library
\item
  \textbf{Native platform feature?} --- use platform primitives
\item
  \textbf{Installed dependency?} --- use what's already there
\item
  \textbf{One line?} --- write one line
\item
  \textbf{Then}: minimum viable implementation
\end{enumerate}

\vspace{0.3em}
\begin{block}{Non-negotiable}
Validation, error handling, security, and accessibility are always kept — ponytail targets complexity, not safety.
\end{block}
\end{frame}

\begin{frame}[fragile]{Ponytail in Practice --- QPUBench Examples}
\phantomsection\label{ponytail-in-practice-qpubench-examples}
Fixes made during QPUBench development under ponytail review:

\begin{longtable}[]{@{}
  >{\raggedright\arraybackslash}p{(\columnwidth - 4\tabcolsep) * \real{0.3333}}
  >{\raggedright\arraybackslash}p{(\columnwidth - 4\tabcolsep) * \real{0.3333}}
  >{\raggedright\arraybackslash}p{(\columnwidth - 4\tabcolsep) * \real{0.3333}}@{}}
\toprule\noalign{}
\begin{minipage}[b]{\linewidth}\raggedright
File
\end{minipage} & \begin{minipage}[b]{\linewidth}\raggedright
Issue
\end{minipage} & \begin{minipage}[b]{\linewidth}\raggedright
Fix
\end{minipage} \\
\midrule\noalign{}
\endhead
\texttt{store.py} \texttt{ParquetStore.query} & \texttt{pd.Series} mask
--- bug + unnecessary & Replaced with
\texttt{df\ =\ df{[}df{[}col{]}\ ==\ val{]}} loop \\
\texttt{stub.py} \texttt{StubMBQCAdapter.spec} &
\texttt{pattern\ =\ None} --- dead variable & Removed \\
\texttt{adapter.py} (qforte) & \texttt{tuple({[}sym,\ tuple(xyz){]})}
--- list wrapped in tuple & Changed to \texttt{(sym,\ tuple(xyz))} \\
\texttt{store.py} \texttt{\_iter} & No return type & Added
\texttt{-\textgreater{}\ Iterator{[}BenchmarkRecord{]}} \\
\texttt{record.py} \texttt{\_utcnow} & Lambda assigned to name (noqa
E731) & Converted to typed \texttt{def} \\
\bottomrule\noalign{}
\end{longtable}

\textbf{Result}: fewer lines, cleaner types, one real bug found.
\end{frame}

\begin{frame}[fragile]{Python Quality Toolchain --- QPUBench}
\phantomsection\label{python-quality-toolchain-qpubench}
\footnotesize

\begin{Shaded}
\begin{Highlighting}[]
\ExtensionTok{ruff}\NormalTok{ check src/ tests/    }\CommentTok{\# lint (line{-}length 100, target py311)}
\ExtensionTok{ruff}\NormalTok{ format src/ tests/   }\CommentTok{\# format}
\ExtensionTok{mypy}\NormalTok{ src/                 }\CommentTok{\# static types (strict = true)}
\ExtensionTok{pytest}\NormalTok{ tests/ }\AttributeTok{{-}q}          \CommentTok{\# 259 tests, SDK{-}dependent ones skip cleanly}
\end{Highlighting}
\end{Shaded}

\begin{longtable}[]{@{}
  >{\raggedright\arraybackslash}p{(\columnwidth - 2\tabcolsep) * \real{0.5000}}
  >{\raggedright\arraybackslash}p{(\columnwidth - 2\tabcolsep) * \real{0.5000}}@{}}
\toprule\noalign{}
\begin{minipage}[b]{\linewidth}\raggedright
Tool
\end{minipage} & \begin{minipage}[b]{\linewidth}\raggedright
Role
\end{minipage} \\
\midrule\noalign{}
\endhead
\textbf{Pydantic v2} & Runtime-validated schemas --- field types
enforced at parse time \\
\textbf{ruff} & Linting + auto-formatting, replaces
flake8/isort/black \\
\textbf{mypy strict} & Full type coverage --- \texttt{Any} only where
unavoidable \\
\textbf{pytest} & Schema-only test suite runs with
\texttt{pip\ install\ .} alone \\
\bottomrule\noalign{}
\end{longtable}

\normalsize
\end{frame}

\begin{frame}[fragile]{Ponytail --- Install \& Use}
\phantomsection\label{ponytail-install-use}
\begin{Shaded}
\begin{Highlighting}[]
\CommentTok{\# Claude Code plugin}
\ExtensionTok{/plugin}\NormalTok{ install ponytail@ponytail}

\CommentTok{\# Review current diff for over{-}engineering}
\ExtensionTok{/ponytail{-}review}

\CommentTok{\# Scan entire repo for unnecessary code}
\ExtensionTok{/ponytail{-}audit}

\CommentTok{\# Collect deferred optimisation shortcuts}
\ExtensionTok{/ponytail{-}debt}
\end{Highlighting}
\end{Shaded}

\textbf{Supported agents}: Claude Code · OpenAI Codex · GitHub Copilot
CLI · Cursor · Windsurf · Gemini CLI · Devin

\vspace{0.3em}

\emph{github.com/DietrichGebert/ponytail}
\end{frame}

\section{Multi-Agent System}\label{multi-agent-system}

\begin{frame}[fragile]{Why Multi-Agent?}
\phantomsection\label{why-multi-agent}
\textbf{Single-agent limitation}: one agent, one context window, one
repo at a time.

\textbf{Multi-agent enables}:

\begin{itemize}
\tightlist
\item
  Parallel work across multiple repos or services
\item
  Specialised agents per domain (schemas · adapters · tests · docs)
\item
  Continuous background tasks (CI watching, PR review, dependency
  updates)
\item
  Self-hosted pipelines with no cloud data exposure
\end{itemize}

\vspace{0.3em}

\begin{verbatim}
User
 └── Orchestrator agent
         ├── Schema agent     (qpubench schemas)
         ├── Adapter agent    (integrations/)
         └── Review agent     (tests + AGENTS.md checks)
\end{verbatim}
\end{frame}

\begin{frame}[fragile]{gitea-mcp Integration}
\phantomsection\label{gitea-mcp-integration}
\textbf{gitea-mcp} connects AI agents directly to Gitea repositories via
MCP

\begin{verbatim}
AI Agent (Claude Code / Cursor / Windsurf)
    └── MCP protocol
            └── gitea-mcp server  (gitea.com/gitea/gitea-mcp)
                    └── Gitea REST API
                            └── repos · issues · PRs · wikis · CI
\end{verbatim}

\begin{itemize}
\tightlist
\item
  Official MCP server from the Gitea project
\item
  Agents can read/write code, open PRs, comment on issues, trigger CI
\item
  Works with Claude Code, Mistral Vibe, Gemini CLI, OpenCode, Devin
\item
  Self-hosted Gitea = full data sovereignty
\end{itemize}
\end{frame}

\begin{frame}[fragile]{Supported Agents \& Platforms}
\phantomsection\label{supported-agents-platforms}
\begin{longtable}[]{@{}llll@{}}
\toprule\noalign{}
Agent & Open-weight & Self-hostable & gitea-mcp \\
\midrule\noalign{}
\endhead
Claude Code (Anthropic) & No & No & Yes \\
Mistral Vibe & Yes & Yes & Yes \\
Gemini CLI (Google) & No & No & Yes \\
GitHub Copilot CLI & No & No & Yes \\
OpenCode & Depends & Yes & Yes \\
Devin & No & No & Yes \\
\bottomrule\noalign{}
\end{longtable}

\vspace{0.3em}

All agents read the same \texttt{AGENTS.md} file --- one set of rules,
any tool.
\end{frame}

\begin{frame}[fragile]{Claude Code + Mistral Vibe}
\phantomsection\label{claude-code-mistral-vibe}
\begin{columns}[T]
\begin{column}{0.5\textwidth}
\textbf{Claude Code} \emph{(Anthropic)}

\begin{itemize}
\tightlist
\item
  CLI-first agentic coding assistant
\item
  \texttt{CLAUDE.md} / \texttt{AGENTS.md} for persistent rules
\item
  Git worktrees for isolated parallel agents
\item
  Per-project memory system across sessions
\item
  Hooks for automated validation
\end{itemize}
\end{column}

\begin{column}{0.5\textwidth}
\textbf{Mistral Vibe} \emph{(Mistral AI)}

\begin{itemize}
\tightlist
\item
  Open-weight models, self-hostable
\item
  MCP-compatible via gitea-mcp
\item
  Suitable for offline / air-gapped environments
\item
  No data leaves your infrastructure
\end{itemize}
\end{column}
\end{columns}

\vspace{0.3em}

Both integrate with the same \texttt{AGENTS.md} instruction file.
\end{frame}

\begin{frame}[fragile]{Context \& Memory Management}
\phantomsection\label{context-memory-management}
\textbf{Per-session context discipline}:

\begin{longtable}[]{@{}
  >{\raggedright\arraybackslash}p{(\columnwidth - 2\tabcolsep) * \real{0.4000}}
  >{\raggedright\arraybackslash}p{(\columnwidth - 2\tabcolsep) * \real{0.6000}}@{}}
\toprule\noalign{}
\begin{minipage}[b]{\linewidth}\raggedright
File
\end{minipage} & \begin{minipage}[b]{\linewidth}\raggedright
Purpose
\end{minipage} \\
\midrule\noalign{}
\endhead
\texttt{AGENTS.md} & Language-agnostic, permanent agent instructions \\
\texttt{CLAUDE.md} & Claude-specific rules, project conventions \\
\texttt{CONTEXT.md} & Current task, status, constraints --- written
before each session \\
\bottomrule\noalign{}
\end{longtable}

\textbf{Workflow}:

\begin{itemize}
\tightlist
\item
  Before starting → write \texttt{CONTEXT.md} with current task +
  constraints
\item
  Each session → \texttt{"Read\ CONTEXT.md\ and\ CLAUDE.md\ first"}
\item
  When switching repos → update \texttt{CONTEXT.md} with what was just
  completed
\item
  Parallel agents → clear boundary definitions + shared
  \texttt{contract.ts}
\end{itemize}
\end{frame}

\begin{frame}[fragile]{Multi-Agent Memory Architecture}
\phantomsection\label{multi-agent-memory-architecture}
\textbf{Four persistent memory types} (survive across sessions):

\begin{longtable}[]{@{}lll@{}}
\toprule\noalign{}
Type & Content & When to save \\
\midrule\noalign{}
\endhead
\texttt{user} & Role, expertise, preferences & First conversation \\
\texttt{feedback} & Corrections + validated approaches & After any
feedback \\
\texttt{project} & Goals, deadlines, decisions & When context changes \\
\texttt{reference} & External system pointers & When a resource is
named \\
\bottomrule\noalign{}
\end{longtable}

\begin{verbatim}
~/.claude/projects/<repo>/memory/
├── MEMORY.md            <- index (always loaded, max 200 lines)
├── user_role.md
├── feedback_testing.md
└── project_schema.md
\end{verbatim}
\end{frame}

\begin{frame}[fragile]{Self-Deployed LLMs}
\phantomsection\label{self-deployed-llms}
\textbf{Abzu} (HuggingFace) --- open-weight models for quantum computing
tasks

\begin{itemize}
\tightlist
\item
  Local deployment via HuggingFace \texttt{transformers} / Ollama / vLLM
\item
  No data leaves your infrastructure
\item
  Suitable for air-gapped laboratory environments
\item
  Combine with self-hosted Gitea + gitea-mcp for a fully sovereign
  pipeline
\end{itemize}

\vspace{0.3em}

\begin{verbatim}
Local LLM (Ollama / vLLM)
    └── MCP client (Claude Code / Cursor)
            └── gitea-mcp  -->  Gitea (self-hosted)
                                    └── qpubench repo
\end{verbatim}

\begin{block}{Full self-hosted stack}
Model · version control · CI · agent memory — nothing leaves your network.
\end{block}
\end{frame}

\section{Summary}\label{summary}

\begin{frame}[fragile]{Key Takeaways}
\phantomsection\label{key-takeaways}
\begin{columns}[T]
\begin{column}{0.5\textwidth}
\textbf{Lean code (ponytail)}

\begin{itemize}
\tightlist
\item
  Decision ladder before every change
\item
  Runtime enforcement \textgreater\textgreater{} documentation
\item
  Remove dead code ruthlessly
\item
  Type hints everywhere
\item
  46\% fewer lines, 22\% fewer tokens
\end{itemize}
\end{column}

\begin{column}{0.5\textwidth}
\textbf{Multi-agent systems}

\begin{itemize}
\tightlist
\item
  gitea-mcp bridges agents to Gitea
\item
  Same \texttt{AGENTS.md} for all tools
\item
  \texttt{CONTEXT.md} per session
\item
  Four memory types persist state
\item
  Full self-hosted stack possible
\end{itemize}
\end{column}
\end{columns}

\vspace{0.3em}

\textbf{Links}: \texttt{github.com/DietrichGebert/ponytail} ·
\texttt{gitea.com/gitea/gitea-mcp} · \texttt{huggingface.co/Abzu}
\end{frame}

\end{document}
